\chapter{总结与展望}

\section{工作总结}

本作品围绕“基于多类格困难假设的抗量子身份基匹配加密通用设计和优化实现”完成了从理论构造、
密码原语到应用入口的完整实现。系统以两层匿名HIBE承载接收身份和认可发送身份，以IBS认证
发送者，并通过短种子混合加密承载不同长度的应用消息。完整密文进一步绑定一次性验签公钥和
全部传输字节，使身份条件、载荷与外层字节认证处于同一处理流程。

在算法层，作品实现LATTE、Lai两种HIBE，以及BNN＋ML-DSA、BNN＋Falcon、Luo、SRGG、
Sageloli五种常规IBS，形成十种端到端组合；在此基础上又接入Wang--PW目标参数矩阵SIS后端。
三组代表分别覆盖低延迟交互、原生模格身份签名和固定身份对连续通信。所有组合复用相同的公共
五算法、对象边界、状态管理、命令行和浏览器入口。

在优化层，作品针对完整运行中出现的主要瓶颈提出三项实现方法：Sageloli结构化采样保留环结构，
把大规模相关扰动转化为逐频率低秩计算；SRGG范围感知编码按字段实际取值范围压缩四类对象；
Lai接收会话复用固定身份对的叶密钥，将重复委派移到准备阶段。三项优化都不改变相应算法的
公开方程，也不要求修改通用IB-ME协议。

在软件层，后端通过稳定接口、能力声明和规范编解码接入。组合预检在建立系统前确认身份宽度、
签名输入、载荷容量、对象格式和运行依赖；固定发钥服务与原子状态保存支持跨进程恢复；CLI和
Web界面完成密钥建立、角色导出、定向加密、匹配解密及错误输入演示。

\section{实验结论}

A组重复加密、解密中位数分别为78.5毫秒和152毫秒，适合交互式演示和定向消息。B组的主要
时间集中在发送者发钥和SRGG签名，重复解密中位数为161毫秒；三档SRGG主状态的紧凑长度为
JSON的18.3\%至19.8\%，说明字段范围对大对象存储具有明显影响。

C组一次完整解密中，Lai委派约占92.5\%。接收会话把该阶段前移后，会话准备约510秒，后续
解密中位数为28.6秒。Sageloli PARAM-I、II均完成目标参数发钥、签名和公开验证，结构化实现
的峰值RSS分别约1.36~GiB和6.19~GiB。Wang--PW完成目标规模建立、身份发钥、独立签名、公开
验签及LATTE组合，其大对象和服务器级资源需求也说明通用接口需要同时容纳轻量与重型后端。

密文组成分析表明，A组主要开销来自Lamport一次性认证，C组主要开销来自Sageloli签名和Lai
胶囊。优化单个对象能否显著改变完整传输量，取决于该对象在整包中的占比。身份错配、字节篡改、
截断、追加、内部签名错误和跨进程恢复测试则覆盖了统一协议的主要功能路径。

\section{后续工作}

第一项后续工作是在Ultra~9~275HX同一机器上完成十种常规组合补测。统一操作系统、Python与
Sage版本、线程配置、身份长度、消息长度和重复次数后，可以形成更严格的组合横向对比，并补充
时间分位数、进程树峰值内存、对象尺寸和完整密文长度。

第二项工作是降低协议级固定通信开销。当前Lamport一次性公钥与签名合计96~KiB，在A组短消息
中占据主要比例。后续可研究更紧凑的一次性认证或直接采用具备所需强不可伪造性质的签名构造，
同时保持完整密文字节绑定和现有错误输入语义。

第三项工作是继续扩展HIBE与IBS后端。接入过程已经固定为接口实现、能力声明、规范编解码、
原语自检和组合测试；后续可以增加其他环格、模格或矩阵格方案，并在同机条件下比较它们在
Setup、发钥、在线收发、持久恢复和通信量方面的差异。

第四项工作是完善面向长期部署的身份与密钥管理，包括带有效期的身份编码、密钥撤销与更新、
硬件隔离的主状态、会话密钥清理和运行审计。浏览器演示已经验证角色操作流程，进一步集成时可
将规范密文字节接入现有消息队列或文件交换系统，并保持密码运算与业务界面的分层。

通过这些工作，系统可以从多算法研究原型继续发展为可复现、可扩展并适应不同资源条件的后量子
身份基匹配加密平台。
